% Автоматаар үүсгэсэн — эх файл: lab04/README.md
\chapter{Лаб 4 — Багтаамжийн хязгаар: ирмэг ба үүл}

\begingroup\scriptsize\begin{longtable}[]{@{}
  >{\raggedright\arraybackslash}p{(\columnwidth - 2\tabcolsep) * \real{0.5000}}
  >{\raggedright\arraybackslash}p{(\columnwidth - 2\tabcolsep) * \real{0.5000}}@{}}
\toprule
\endhead
\bottomrule
\endlastfoot
\textbf{7 хоног} & IX \\
\textbf{Хугацаа} & 4 цаг \\
\textbf{Суралцахуйн үр дүн} & ҮД2 (өргөтгөх боломжтой backend), ҮД4, ҮД6, ҮД8 \\
\textbf{Үнэлгээ} & \textbf{10 минутын багийн ҮЗҮҮЛЭН} (бичгийн тайлан биш), 10 оноо \\
\textbf{Гол хэмжилт} & Аль нөөц ХАМГИЙН ТҮРҮҮНД ханав --- RAM, CPU, 100 Mbit уплинк эсвэл microSD \\
\end{longtable}\endgroup

\hypertarget{ux437ux43eux440ux438ux43bux433ux43e}{%
\section{1. Зорилго}\label{ux437ux43eux440ux438ux43bux433ux43e}}

Энэ бол хичээлийн \textbf{гол лаборатори}. Та нэг брокер биш, \textbf{хоёр брокерыг} ачаалж, тэдгээрийн хязгаарыг харьцуулна:

\begin{itemize}
\tightlist
\item
  {[}Pi{]} Raspberry Pi 3B дээрх \textbf{mosquitto} --- 1 GB RAM, 4×Cortex-A53, 100 Mbit уплинк (USB 2.0), microSD
\item
  {[}ЗК{]} Зөөврийн компьютер дээрх \textbf{EMQX} --- олон GB RAM, олон цөм, SSD
\end{itemize}

Гурван асуултад тоон хариулт өгнө:

\begin{enumerate}
\def\labelenumi{\arabic{enumi}.}
\tightlist
\item
  Хоёр брокер тус бүр хэдэн холболт, хэдэн мессеж/сек даах вэ?
\item
  \textbf{Аль нөөц түрүүлж ханав?} --- RAM, CPU, уплинк, эсвэл microSD. Хариултаа \textbf{нотол}, зүгээр битгий бич.
\item
  Мессежийг хаана боловсруулах вэ --- брокер дээр (EMQX rule engine) үү, тусдаа хэрэглэгч дээр үү?
\end{enumerate}

\begin{quote}
\textbf{Хамгийн чухал зарчим:} "Pi 3B 200 холболт даасан" гэдэг нь хариулт биш. "200 холболт дээр сул RAM 118 MiB болж swap \texttt{si/so} тэгээс өссөн, харин CPU 140\% буюу 4 цөмийн 35\% л байсан --- тиймээс хязгаар нь \textbf{санах ой}" гэдэг нь хариулт. Үзүүлэнгийн оноо энэ ялгаан дээр тогтоно.
\end{quote}

\begin{quote}
\textbf{Анхаар:} энэ лабораторийг \textbf{үзүүлэнгээр} үнэлнэ. Багийн гурван гишүүн бүгд ярина. Слайд шаардлагагүй --- ажиллаж буй график ба тоо хангалттай.
\end{quote}

\hypertarget{ux443ux440ux44cux434ux447ux438ux43bux441ux430ux43d-ux43dux4e9ux445ux446ux4e9ux43b}{%
\section{2. Урьдчилсан нөхцөл}\label{ux443ux440ux44cux434ux447ux438ux43bux441ux430ux43d-ux43dux4e9ux445ux446ux4e9ux43b}}

\begin{itemize}
\tightlist
\item
  Лаб 1--3 дууссан. \textbf{Лаб 1-ийн Хүснэгт 1.5 (ирмэгийн ачааллын хариу үйлдэл) ба Лаб 3-ын Хүснэгт 3.6 (ачааллын хэмжээний нөлөө) гар дор байх ёстой} --- өнөөдөр тэдгээрээс таамаг гаргана.
\item
  {[}ЗК{]} Файлын дескрипторын хязгаар өргөжсөн, дүрс татагдсан:
\end{itemize}

\begin{Shaded}
\begin{Highlighting}[]
\BuiltInTok{ulimit} \AttributeTok{{-}n}                     \CommentTok{\# 1024 бол хангалтгүй}
\BuiltInTok{ulimit} \AttributeTok{{-}n}\NormalTok{ 65535               }\CommentTok{\# Linux/macOS; Windows дээр WSL2 дотор}
\ExtensionTok{docker}\NormalTok{ pull emqx/emqtt{-}bench:0.4.20}
\end{Highlighting}
\end{Shaded}

\begin{itemize}
\tightlist
\item
  {[}Pi{]} Pi бэлэн: \texttt{cd\ \textasciitilde{}/cnc302/edge\ \&\&\ make\ check\ \&\&\ make\ mem} → бүгд OK, throttle \texttt{0x0}
\item
  {[}ЗК{]} EMQX-ийн нууц үг терминалд бэлэн:
\end{itemize}

\begin{Shaded}
\begin{Highlighting}[]
\BuiltInTok{export} \VariableTok{EMQX\_DASHBOARD\_PASSWORD}\OperatorTok{=}\StringTok{\textquotesingle{}\textless{}stack/.env дэх утга\textgreater{}\textquotesingle{}}
\end{Highlighting}
\end{Shaded}

\begin{itemize}
\tightlist
\item
  {[}Pi{]} Pi дээр терминал бүрд: \texttt{set\ -a\ \&\&\ source\ \textasciitilde{}/cnc302/edge/.env\ \&\&\ set\ +a}
\end{itemize}

\hypertarget{ux43eux43dux43eux43bux44bux43d-ux441ux430ux43dux443ux443ux43bux433ux430}{%
\section{3. Онолын сануулга}\label{ux43eux43dux43eux43bux44bux43d-ux441ux430ux43dux443ux443ux43bux433ux430}}

\hypertarget{ux445ux430ux43dux430ux43bux442-saturation-ux433ux443ux440ux432ux430ux43d-ux4afux435-ux448ux430ux442ux442ux430ux439}{%
\subsection{Ханалт (saturation) гурван үе шаттай}\label{ux445ux430ux43dux430ux43bux442-saturation-ux433ux443ux440ux432ux430ux43d-ux4afux435-ux448ux430ux442ux442ux430ux439}}

Ачаалал нэмэгдэхэд систем эхлээд \textbf{шугаман} өснө, дараа нь \textbf{тэгширнэ}, эцэст нь \textbf{буурна}. Сүүлийн үе шат хамгийн аюултай: систем ажиллаж байгаа мэт харагдаад бодит хүчин чадал нь буурсан байдаг. Түүнийг олж харах нь энэ лабораторийн зорилго.

\hypertarget{ux434ux4e9ux440ux432ux4e9ux43d-ux431ux43eux43bux43eux43cux436ux438ux442-ux445ux44fux437ux433ux430ux430ux440ux43bux430ux433ux447-ux431ux430-ux442ux44dux434ux433ux44dux44dux440ux438ux439ux43d-ux433ux430ux440ux44bux43d-ux4afux441ux44dux433}{%
\subsection{Дөрвөн боломжит хязгаарлагч ба тэдгээрийн гарын үсэг}\label{ux434ux4e9ux440ux432ux4e9ux43d-ux431ux43eux43bux43eux43cux436ux438ux442-ux445ux44fux437ux433ux430ux430ux440ux43bux430ux433ux447-ux431ux430-ux442ux44dux434ux433ux44dux44dux440ux438ux439ux43d-ux433ux430ux440ux44bux43d-ux4afux441ux44dux433}}

\begingroup\scriptsize\begin{longtable}[]{@{}
  >{\raggedright\arraybackslash}p{(\columnwidth - 4\tabcolsep) * \real{0.3333}}
  >{\raggedright\arraybackslash}p{(\columnwidth - 4\tabcolsep) * \real{0.3333}}
  >{\raggedright\arraybackslash}p{(\columnwidth - 4\tabcolsep) * \real{0.3333}}@{}}
\toprule
\begin{minipage}[b]{\linewidth}\raggedright
Хязгаарлагч
\end{minipage} & \begin{minipage}[b]{\linewidth}\raggedright
Ямар тоогоор илэрнэ
\end{minipage} & \begin{minipage}[b]{\linewidth}\raggedright
Хаанаас харах
\end{minipage} \\
\midrule
\endhead
\bottomrule
\endlastfoot
\textbf{RAM} & сул RAM буурч 150 MiB-ээс доош, swap \texttt{si/so} \textgreater{} 0 & \texttt{measure\_stack.sh}, \texttt{vmstat\ 1} \\
\textbf{CPU} & нийт CPU 380--400\% (4 цөм × 100\%) & \texttt{docker\ stats}, \texttt{cpu\_\allowbreak{}total\_\allowbreak{}pct} \\
\textbf{Уплинк} & Mbit/с 90--95-д ойртоно, саатал огцом өснө & \texttt{docker\ stats} NET I/O, тооцоо \\
\textbf{microSD} & I/O хүлээлт (\texttt{wa}) өснө, mosquitto-гийн persistence удаашрана & \texttt{vmstat\ 1}, \texttt{iostat} \\
\end{longtable}\endgroup

Pi 3B дээр \textbf{ихэвчлэн RAM түрүүлдэг, CPU биш} --- гэхдээ энэ бол таамаг. Өнөөдөр үүнийг батлах эсвэл няцаах нь та.

\hypertarget{ux443ux43fux43bux438ux43dux43aux438ux439ux43d-ux430ux440ux438ux444ux43cux435ux442ux438ux43a-ux445ux44dux43cux436ux438ux445ux44dux44dux441ux44dux44d-ux4e9ux43cux43dux4e9-ux442ux43eux43eux446ux43eux43eux43b}{%
\subsection{Уплинкийн арифметик --- хэмжихээсээ ӨМНӨ тооцоол}\label{ux443ux43fux43bux438ux43dux43aux438ux439ux43d-ux430ux440ux438ux444ux43cux435ux442ux438ux43a-ux445ux44dux43cux436ux438ux445ux44dux44dux441ux44dux44d-ux4e9ux43cux43dux4e9-ux442ux43eux43eux446ux43eux43eux43b}}

\begin{verbatim}
Bitrate (Mbit/с) = N_холболт × R_мсж/с × S_байт × 8 / 1 000 000
\end{verbatim}

\texttt{loadtest.sh} нь \texttt{pub} горимд яг үүнийг өөрөө хэвлэдэг. 100 Mbit холбоос дээр бодит хүрэх дээд хурд \textbf{\textasciitilde90--95 Mbit/с} бөгөөд Ethernet нь USB 2.0-тэй зурвасаа хуваана.

\begin{quote}
\textbf{Урхи:} \texttt{loadtest.sh}-ийн анхдагч сэдэв нь \texttt{cnc302/\allowbreak{}bench/\allowbreak{}\%i/\allowbreak{}telemetry}. Гүүр нь зөвхөн \texttt{cnc302/\textless{}SITE\textgreater{}/\#}-ийг үүл рүү дамжуулдаг тул \textbf{анхдагч ачаалал уплинкийг огт хөндөхгүй} --- зөвхөн Pi-гийн mosquitto-г ачаална. Уплинкийг ачаалах бол \texttt{TOPIC}-ыг гүүрээр дамжих угтвартай болгох ёстой (Алхам 3.3).
\end{quote}

\hypertarget{ux445ux430ux430ux43dux430-ux431ux43eux43bux43eux432ux441ux440ux443ux443ux43bux430ux445-ux432ux44d}{%
\subsection{Хаана боловсруулах вэ}\label{ux445ux430ux430ux43dux430-ux431ux43eux43bux43eux432ux441ux440ux443ux443ux43bux430ux445-ux432ux44d}}

\begingroup\scriptsize\begin{longtable}[]{@{}
  >{\raggedright\arraybackslash}p{(\columnwidth - 4\tabcolsep) * \real{0.3333}}
  >{\raggedright\arraybackslash}p{(\columnwidth - 4\tabcolsep) * \real{0.3333}}
  >{\raggedright\arraybackslash}p{(\columnwidth - 4\tabcolsep) * \real{0.3333}}@{}}
\toprule
\begin{minipage}[b]{\linewidth}\raggedright
Байрлал
\end{minipage} & \begin{minipage}[b]{\linewidth}\raggedright
Давуу тал
\end{minipage} & \begin{minipage}[b]{\linewidth}\raggedright
Сул тал
\end{minipage} \\
\midrule
\endhead
\bottomrule
\endlastfoot
Брокер дээр (EMQX rule engine) & хамгийн бага саатал, нэмэлт үсрэлт үгүй & брокерын CPU иднэ, SQL-ээр хязгаарлагдмал \\
Тусдаа хэрэглэгч (Node-RED, Python) & бүрэн уян хатан, тусад нь өргөтгөнө & нэмэлт үсрэлт, нэмэлт үйлчилгээ \\
Ирмэг дээр (агент) & уплинк хэмнэнэ, холбоос тасарсан ч ажиллана & Pi-гийн CPU/RAM хязгаартай \\
\end{longtable}\endgroup

\hypertarget{ux430ux43bux445ux43cux443ux443ux434}{%
\section{4. Алхмууд}\label{ux430ux43bux445ux43cux443ux443ux434}}

\hypertarget{ux430ux43bux445ux430ux43c-1-ux441ux443ux443ux440ux44c-ux442ux4e9ux43bux4e9ux432-ux445ux43eux451ux440-ux445ux43eux441ux442-ux434ux44dux44dux440-20-ux43cux438ux43d-piux437ux43a}{%
\subsection{Алхам 1 --- Суурь төлөв, хоёр хост дээр (20 мин) {[}Pi{]}{[}ЗК{]}}\label{ux430ux43bux445ux430ux43c-1-ux441ux443ux443ux440ux44c-ux442ux4e9ux43bux4e9ux432-ux445ux43eux451ux440-ux445ux43eux441ux442-ux434ux44dux44dux440-20-ux43cux438ux43d-piux437ux43a}}

Ачаалал өгөхийн өмнөх төлөв. Дараагийн бүх тоо үүнтэй харьцуулагдана.

\begin{Shaded}
\begin{Highlighting}[]
\CommentTok{\# [Pi] Pi дээр}
\ExtensionTok{pkill} \AttributeTok{{-}f}\NormalTok{ vscode{-}server                       }\CommentTok{\# 150–250 MiB чөлөөлнө}
\BuiltInTok{cd}\NormalTok{ \textasciitilde{}/cnc302 }\KeywordTok{\&\&} \ExtensionTok{./tools/measure\_stack.sh} \AttributeTok{{-}{-}role}\NormalTok{ edge}
\ExtensionTok{mosquitto\_sub} \AttributeTok{{-}h}\NormalTok{ localhost }\AttributeTok{{-}t} \StringTok{\textquotesingle{}$SYS/broker/clients/connected\textquotesingle{}} \AttributeTok{{-}C}\NormalTok{ 1}
\ExtensionTok{mosquitto\_sub} \AttributeTok{{-}h}\NormalTok{ localhost }\AttributeTok{{-}t} \StringTok{\textquotesingle{}$SYS/broker/messages/stored\textquotesingle{}} \AttributeTok{{-}C}\NormalTok{ 1}
\CommentTok{\# [ЗК] компьютер дээр}
\BuiltInTok{cd}\NormalTok{ \textasciitilde{}/cnc302 }\KeywordTok{\&\&} \ExtensionTok{./tools/measure\_stack.sh} \AttributeTok{{-}{-}role}\NormalTok{ cloud}
\end{Highlighting}
\end{Shaded}

\hypertarget{ux445ux4afux441ux43dux44dux433ux442-4.1-ux441ux443ux443ux440ux44c-ux442ux4e9ux43bux4e9ux432}{%
\subsubsection{Хүснэгт 4.1 --- Суурь төлөв}\label{ux445ux4afux441ux43dux44dux433ux442-4.1-ux441ux443ux443ux440ux44c-ux442ux4e9ux43bux4e9ux432}}

\begingroup\scriptsize\begin{longtable}[]{@{}lll@{}}
\toprule
Хэмжигдэхүүн & {[}Pi{]} ирмэг (Pi 3B) & {[}ЗК{]} үүл (компьютер) \\
\midrule
\endhead
\bottomrule
\endlastfoot
Сул RAM (MiB) & & \\
Swap ашигласан (MiB) & & \\
CPU \% (сул үед) & & \\
Температур (°C) / \texttt{get\_throttled} & / & --- \\
Брокерын одоогийн холболт & & \\
Брокерын RAM (MiB) & & \\
\end{longtable}\endgroup

\hypertarget{ux430ux43bux445ux430ux43c-2-ux442ux430ux430ux43cux430ux433-ux445ux44fux437ux433ux430ux430ux440ux44bux433-ux445ux44dux43cux436ux438ux445ux44dux44dux441ux44dux44d-ux4e9ux43cux43dux4e9-ux442ux43eux43eux446ux43eux43eux43b-25-ux43cux438ux43d-ux437ux43a}{%
\subsection{Алхам 2 --- Таамаг: хязгаарыг ХЭМЖИХЭЭСЭЭ ӨМНӨ тооцоол (25 мин) {[}ЗК{]}}\label{ux430ux43bux445ux430ux43c-2-ux442ux430ux430ux43cux430ux433-ux445ux44fux437ux433ux430ux430ux440ux44bux433-ux445ux44dux43cux436ux438ux445ux44dux44dux441ux44dux44d-ux4e9ux43cux43dux4e9-ux442ux43eux43eux446ux43eux43eux43b-25-ux43cux438ux43d-ux437ux43a}}

Багаараа цаасан дээр тооцоолж, \textbf{дараа нь битгий засаарай}. Буруу таамаг бол сонирхолтой үр дүн; засварласан таамаг бол оноогүй.

\textbf{2.1 Уплинкийн таазыг тооцоол} (200 байт ачаалал, 1 мсж/с нэг холболтод):

\begin{verbatim}
1 мессеж = 200 Б × 8 = 1 600 бит
100 Mbit/с ÷ 1 600 бит = ______ мсж/с   ← онолын дээд хязгаар
Ачааллын шатууд 10…500 холболт × 1 мсж/с = 10…500 мсж/с
→ Уплинк ханах уу? ______  (Хэдэн % ашиглагдах вэ: ______)
\end{verbatim}

\textbf{2.2 Санах ойн таазыг тооцоол.} Лаб 1-ийн Хүснэгт 1.5-аас нэг холболтод ногдох RAM-ыг гарга:

\begin{verbatim}
ΔRAM (MiB) / Δхолболт = ______ KiB/холболт
Сул RAM (Хүснэгт 4.1) − 150 MiB нөөц = ______ MiB ашиглаж болно
→ RAM ханах холболтын тоо ≈ ______
\end{verbatim}

\textbf{2.3 CPU-гийн таазыг тооцоол.} Лаб 1-д хэмжсэн нэг мессежид ногдох CPU-гаас:

\begin{verbatim}
4 цөм × 100% = 400% боломжтой;  70% ашиглалтаар = 280%
→ CPU ханах мессежийн хурд ≈ ______ мсж/с
\end{verbatim}

\hypertarget{ux445ux4afux441ux43dux44dux433ux442-4.2-ux442ux430ux430ux43cux430ux433-ux445ux44dux43cux436ux438ux445ux44dux44dux441-ux4e9ux43cux43dux4e9-ux431ux4e9ux433ux43bux4e9ux43dux4e9}{%
\subsubsection{Хүснэгт 4.2 --- Таамаг (хэмжихээс ӨМНӨ бөглөнө)}\label{ux445ux4afux441ux43dux44dux433ux442-4.2-ux442ux430ux430ux43cux430ux433-ux445ux44dux43cux436ux438ux445ux44dux44dux441-ux4e9ux43cux43dux4e9-ux431ux4e9ux433ux43bux4e9ux43dux4e9}}

\begingroup\scriptsize\begin{longtable}[]{@{}lll@{}}
\toprule
Нөөц & Ханах цэг (холболт эсвэл мсж/с) & Тооцооны үндэс \\
\midrule
\endhead
\bottomrule
\endlastfoot
RAM & & \\
CPU & & \\
100 Mbit уплинк & & \\
microSD (persistence) & & \\
\textbf{Бидний таамаг: ХАМГИЙН ТҮРҮҮНД ханах нөөц} & & \\
\end{longtable}\endgroup

\hypertarget{ux430ux43bux445ux430ux43c-3-ux438ux440ux43cux44dux433ux438ux439ux43d-ux430ux447ux430ux430ux43bux43bux44bux43d-ux442ux435ux441ux442-60-ux43cux438ux43d-ux437ux43api}{%
\subsection{Алхам 3 --- Ирмэгийн ачааллын тест (60 мин) {[}ЗК{]}{[}Pi{]}}\label{ux430ux43bux445ux430ux43c-3-ux438ux440ux43cux44dux433ux438ux439ux43d-ux430ux447ux430ux430ux43bux43bux44bux43d-ux442ux435ux441ux442-60-ux43cux438ux43d-ux437ux43api}}

\textbf{3.1 Дасгал ажиллуулалт} --- бүх зүйл холбогдож байгааг батлах:

\begin{Shaded}
\begin{Highlighting}[]
\CommentTok{\# [ЗК] компьютер дээрээс Pi рүү (ачаалал ҮРГЭЛЖ өөр хостоос ирнэ)}
\BuiltInTok{cd}\NormalTok{ \textasciitilde{}/cnc302 }\KeywordTok{\&\&} \FunctionTok{bash}\NormalTok{ lab04/loadtest.sh }\AttributeTok{{-}{-}target}\NormalTok{ edge pub }\OperatorTok{\textless{}}\NormalTok{PI{-}IP}\OperatorTok{\textgreater{}}\NormalTok{ 10 1}
\end{Highlighting}
\end{Shaded}

30 секунд ажиллуулаад Ctrl+C. {[}Pi{]} Pi дээр холболт харагдах ёстой:

\begin{Shaded}
\begin{Highlighting}[]
\ExtensionTok{mosquitto\_sub} \AttributeTok{{-}h}\NormalTok{ localhost }\AttributeTok{{-}t} \StringTok{\textquotesingle{}$SYS/broker/clients/connected\textquotesingle{}} \AttributeTok{{-}C}\NormalTok{ 1   }\CommentTok{\# ≈ 10}
\end{Highlighting}
\end{Shaded}

\textbf{3.2 Бүтэн шат дараалсан тест.} Гурван терминалыг \textbf{зэрэг} ажиллуулна:

\begin{Shaded}
\begin{Highlighting}[]
\CommentTok{\# [Pi] терминал 1 — систем хэмжилт}
\BuiltInTok{cd}\NormalTok{ \textasciitilde{}/cnc302 }\KeywordTok{\&\&} \ExtensionTok{./tools/measure\_stack.sh} \AttributeTok{{-}{-}role}\NormalTok{ edge }\AttributeTok{{-}{-}watch}\NormalTok{ 900 5}

\CommentTok{\# [Pi] терминал 2 — үзүүлэлт цуглуулагч (CSV{-}д \textasciigrave{}target=edge\textasciigrave{} бичигдэнэ)}
\BuiltInTok{cd}\NormalTok{ \textasciitilde{}/cnc302 }\KeywordTok{\&\&} \ExtensionTok{python3}\NormalTok{ lab04/collect\_metrics.py }\AttributeTok{{-}{-}target}\NormalTok{ edge }\AttributeTok{{-}{-}label}\NormalTok{ ramp }\DataTypeTok{\textbackslash{}}
    \AttributeTok{{-}{-}seconds}\NormalTok{ 900 }\AttributeTok{{-}{-}emqx} \StringTok{"http://}\VariableTok{$CLOUD\_HOST}\StringTok{:18083"}

\CommentTok{\# [Pi] терминал 3 — swap{-}ын хяналт. si/so \textgreater{} 0 бол ХЯЗГААР ОЛДЛОО}
\FunctionTok{vmstat}\NormalTok{ 1}

\CommentTok{\# [ЗК] ачаалал үүсгэгч — шатууд: 10 25 50 100 200 350 500}
\FunctionTok{bash}\NormalTok{ lab04/loadtest.sh }\AttributeTok{{-}{-}target}\NormalTok{ edge ramp }\OperatorTok{\textless{}}\NormalTok{PI{-}IP}\OperatorTok{\textgreater{}}
\end{Highlighting}
\end{Shaded}

\begin{quote}
⚠ \textbf{\texttt{vmstat\ 1}-ийн \texttt{si}/\texttt{so} багана тэгээс өссөн мөчид зогсоо.} Тэр бол хязгаар. Цааш шахах нь Pi-г царцааж, microSD-г элээнэ. Энэ нь алдаа биш --- \textbf{энэ бол таны хэмжилт}.
\end{quote}

\begin{quote}
\textbf{Санамж:} ирмэгийн ачаалал mosquitto руу очиж байгаа тул CSV дэх \texttt{emqx\_connections} багана \textbf{үүлний} EMQX-ийг харуулна, ачааллыг биш. Ирмэгийн холболтын тоог шат бүрд \texttt{\$SYS/\allowbreak{}broker/\allowbreak{}clients/\allowbreak{}connected}-оос \textbf{гараар} бичиж ав. Энэ өөрөө сургамж: хэмжилтийн хэрэгсэл юуг хэмжиж байгааг үргэлж мэдэж бай.
\end{quote}

Шат бүрд \texttt{emqtt-bench}-ийн гаралт (илгээсэн хурд) ба Pi-гийн үзүүлэлтийг бичнэ:

\hypertarget{ux445ux4afux441ux43dux44dux433ux442-4.3-ux438ux440ux43cux44dux433-pi-3b-mosquitto-ux448ux430ux442ux443ux443ux434}{%
\subsubsection{Хүснэгт 4.3 --- Ирмэг (Pi 3B, mosquitto): шатууд}\label{ux445ux4afux441ux43dux44dux433ux442-4.3-ux438ux440ux43cux44dux433-pi-3b-mosquitto-ux448ux430ux442ux443ux443ux434}}

\begingroup\scriptsize\begin{longtable}[]{@{}
  >{\raggedright\arraybackslash}p{(\columnwidth - 18\tabcolsep) * \real{0.1000}}
  >{\raggedright\arraybackslash}p{(\columnwidth - 18\tabcolsep) * \real{0.1000}}
  >{\raggedright\arraybackslash}p{(\columnwidth - 18\tabcolsep) * \real{0.1000}}
  >{\raggedright\arraybackslash}p{(\columnwidth - 18\tabcolsep) * \real{0.1000}}
  >{\raggedright\arraybackslash}p{(\columnwidth - 18\tabcolsep) * \real{0.1000}}
  >{\raggedright\arraybackslash}p{(\columnwidth - 18\tabcolsep) * \real{0.1000}}
  >{\raggedright\arraybackslash}p{(\columnwidth - 18\tabcolsep) * \real{0.1000}}
  >{\raggedright\arraybackslash}p{(\columnwidth - 18\tabcolsep) * \real{0.1000}}
  >{\raggedright\arraybackslash}p{(\columnwidth - 18\tabcolsep) * \real{0.1000}}
  >{\raggedright\arraybackslash}p{(\columnwidth - 18\tabcolsep) * \real{0.1000}}@{}}
\toprule
\begin{minipage}[b]{\linewidth}\raggedright
Холболт
\end{minipage} & \begin{minipage}[b]{\linewidth}\raggedright
Зорилтот мсж/с
\end{minipage} & \begin{minipage}[b]{\linewidth}\raggedright
Бодит илгээсэн
\end{minipage} & \begin{minipage}[b]{\linewidth}\raggedright
mosquitto холболт (\texttt{\$SYS})
\end{minipage} & \begin{minipage}[b]{\linewidth}\raggedright
CPU \% (4 цөм = 400)
\end{minipage} & \begin{minipage}[b]{\linewidth}\raggedright
Сул RAM (MiB)
\end{minipage} & \begin{minipage}[b]{\linewidth}\raggedright
Swap \texttt{si/so}
\end{minipage} & \begin{minipage}[b]{\linewidth}\raggedright
Уплинк Mbit/с
\end{minipage} & \begin{minipage}[b]{\linewidth}\raggedright
°C
\end{minipage} & \begin{minipage}[b]{\linewidth}\raggedright
\texttt{throttled}
\end{minipage} \\
\midrule
\endhead
\bottomrule
\endlastfoot
10 & 10 & & & & & & & & \\
25 & 25 & & & & & & & & \\
50 & 50 & & & & & & & & \\
100 & 100 & & & & & & & & \\
200 & 200 & & & & & & & & \\
350 & 350 & & & & & & & & \\
500 & 500 & & & & & & & & \\
\end{longtable}\endgroup

\textbf{Хэзээ зогсоох вэ:} сул RAM \textless{} 150 MiB, \texttt{si/so} \textgreater{} 0, \texttt{throttled} ≠ \texttt{0x0}, брокер шинэ холболт авахаа больсон (скрипт өөрөө илрүүлж зогсоно), эсвэл хүлээн авалт илгээсний 95\%-иас доош.

\textbf{3.3 Уплинкийг зориудаар ачаалах.} Дээрх шатууд уплинкийг \textbf{хөндөөгүй} (яагаад? --- §3-ын урхи). Одоо сэдвийг гүүрээр дамжих угтвартай болгож, ачааллыг том болгоно:

\begin{Shaded}
\begin{Highlighting}[]
\CommentTok{\# [ЗК] компьютер дээрээс. 200 холболт × 5 мсж/с × 2000 Б ≈ 16 Mbit/с}
\VariableTok{TOPIC}\OperatorTok{=}\StringTok{"cnc302/}\VariableTok{$SITE}\StringTok{/}\VariableTok{$AREA}\StringTok{/}\VariableTok{$LINE}\StringTok{/bench{-}\%i/telemetry"} \VariableTok{PAYLOAD\_SIZE}\OperatorTok{=}\NormalTok{2000 }\DataTypeTok{\textbackslash{}}
  \FunctionTok{bash}\NormalTok{ lab04/loadtest.sh }\AttributeTok{{-}{-}target}\NormalTok{ edge pub }\OperatorTok{\textless{}}\NormalTok{PI{-}IP}\OperatorTok{\textgreater{}}\NormalTok{ 200 5}
\end{Highlighting}
\end{Shaded}

{[}Pi{]} зэрэг ажиглана (\texttt{docker\ stats}-ийн NET I/O өсөлт ба Хүснэгт 4.3-ын Mbit баганатай харьцуул):

\begin{Shaded}
\begin{Highlighting}[]
\ExtensionTok{docker}\NormalTok{ stats }\AttributeTok{{-}{-}no{-}stream} \AttributeTok{{-}{-}format} \StringTok{\textquotesingle{}table \{\{.Name\}\}\textbackslash{}t\{\{.NetIO\}\}\textbackslash{}t\{\{.MemUsage\}\}\textbackslash{}t\{\{.CPUPerc\}\}\textquotesingle{}}
\end{Highlighting}
\end{Shaded}

Ачааллыг 5 → 10 мсж/с болгож давт. Уплинкийн ашиглалт 90 Mbit/с рүү ойртож байна уу, эсвэл түүнээс өмнө RAM/CPU ханаж байна уу?

\hypertarget{ux445ux4afux441ux43dux44dux433ux442-4.4-ux443ux43fux43bux438ux43dux43aux438ux439ux43d-ux430ux447ux430ux430ux43bux430ux43b-ux441ux44dux434ux44dux432-ux433ux4afux4afux440ux44dux44dux440-ux434ux430ux43cux436ux438ux43dux430}{%
\subsubsection{Хүснэгт 4.4 --- Уплинкийн ачаалал (сэдэв гүүрээр дамжина)}\label{ux445ux4afux441ux43dux44dux433ux442-4.4-ux443ux43fux43bux438ux43dux43aux438ux439ux43d-ux430ux447ux430ux430ux43bux430ux43b-ux441ux44dux434ux44dux432-ux433ux4afux4afux440ux44dux44dux440-ux434ux430ux43cux436ux438ux43dux430}}

\begingroup\scriptsize\begin{longtable}[]{@{}
  >{\raggedright\arraybackslash}p{(\columnwidth - 10\tabcolsep) * \real{0.1667}}
  >{\raggedright\arraybackslash}p{(\columnwidth - 10\tabcolsep) * \real{0.1667}}
  >{\raggedright\arraybackslash}p{(\columnwidth - 10\tabcolsep) * \real{0.1667}}
  >{\raggedright\arraybackslash}p{(\columnwidth - 10\tabcolsep) * \real{0.1667}}
  >{\raggedright\arraybackslash}p{(\columnwidth - 10\tabcolsep) * \real{0.1667}}
  >{\raggedright\arraybackslash}p{(\columnwidth - 10\tabcolsep) * \real{0.1667}}@{}}
\toprule
\begin{minipage}[b]{\linewidth}\raggedright
Холболт × мсж/с × Б
\end{minipage} & \begin{minipage}[b]{\linewidth}\raggedright
Тооцоолсон Mbit/с
\end{minipage} & \begin{minipage}[b]{\linewidth}\raggedright
Хэмжсэн Mbit/с
\end{minipage} & \begin{minipage}[b]{\linewidth}\raggedright
Сул RAM (MiB)
\end{minipage} & \begin{minipage}[b]{\linewidth}\raggedright
CPU \%
\end{minipage} & \begin{minipage}[b]{\linewidth}\raggedright
Гүүрээр үүлэнд хүрсэн үү
\end{minipage} \\
\midrule
\endhead
\bottomrule
\endlastfoot
200 × 5 × 2000 & & & & & \\
200 × 10 × 2000 & & & & & \\
ule{1.5cm}{0.4pt} & & & & & \\
\end{longtable}\endgroup

\hypertarget{ux430ux43bux445ux430ux43c-4-ux4afux4afux43bux43dux438ux439-ux430ux447ux430ux430ux43bux43bux44bux43d-ux442ux435ux441ux442-45-ux43cux438ux43d-ux437ux43a}{%
\subsection{Алхам 4 --- Үүлний ачааллын тест (45 мин) {[}ЗК{]}}\label{ux430ux43bux445ux430ux43c-4-ux4afux4afux43bux43dux438ux439-ux430ux447ux430ux430ux43bux43bux44bux43d-ux442ux435ux441ux442-45-ux43cux438ux43d-ux437ux43a}}

Ижил аргачлал, өөр зорилт. Шатууд \textbf{50 100 250 500 1000 2000}.

\begin{Shaded}
\begin{Highlighting}[]
\CommentTok{\# [ЗК] терминал 1 — цуглуулагч (CSV{-}д \textasciigrave{}target=cloud\textasciigrave{})}
\BuiltInTok{cd}\NormalTok{ \textasciitilde{}/cnc302 }\KeywordTok{\&\&} \ExtensionTok{python3}\NormalTok{ lab04/collect\_metrics.py }\AttributeTok{{-}{-}target}\NormalTok{ cloud }\AttributeTok{{-}{-}label}\NormalTok{ ramp }\AttributeTok{{-}{-}seconds}\NormalTok{ 900}

\CommentTok{\# [ЗК] терминал 2 — ачаалал}
\FunctionTok{bash}\NormalTok{ lab04/loadtest.sh }\AttributeTok{{-}{-}target}\NormalTok{ cloud ramp localhost}
\end{Highlighting}
\end{Shaded}

\begin{quote}
Ачаалал үүсгэгч ба брокер нэг машин дээр байвал \textbf{компьютер өөрөө хязгаарлагч} болж мэднэ. Хоёр дахь машинаас (эсвэл багийн нөгөө гишүүний компьютерээс) ачаалж чадвал илүү зөв. Аль нөхцөлд хэмжсэнээ үзүүлэнд заавал хэл.
\end{quote}

\hypertarget{ux445ux4afux441ux43dux44dux433ux442-4.5-ux4afux4afux43b-ux43aux43eux43cux43fux44cux44eux442ux435ux440-emqx-ux448ux430ux442ux443ux443ux434}{%
\subsubsection{Хүснэгт 4.5 --- Үүл (компьютер, EMQX): шатууд}\label{ux445ux4afux441ux43dux44dux433ux442-4.5-ux4afux4afux43b-ux43aux43eux43cux43fux44cux44eux442ux435ux440-emqx-ux448ux430ux442ux443ux443ux434}}

\begingroup\scriptsize\begin{longtable}[]{@{}
  >{\raggedright\arraybackslash}p{(\columnwidth - 12\tabcolsep) * \real{0.1429}}
  >{\raggedright\arraybackslash}p{(\columnwidth - 12\tabcolsep) * \real{0.1429}}
  >{\raggedright\arraybackslash}p{(\columnwidth - 12\tabcolsep) * \real{0.1429}}
  >{\raggedright\arraybackslash}p{(\columnwidth - 12\tabcolsep) * \real{0.1429}}
  >{\raggedright\arraybackslash}p{(\columnwidth - 12\tabcolsep) * \real{0.1429}}
  >{\raggedright\arraybackslash}p{(\columnwidth - 12\tabcolsep) * \real{0.1429}}
  >{\raggedright\arraybackslash}p{(\columnwidth - 12\tabcolsep) * \real{0.1429}}@{}}
\toprule
\begin{minipage}[b]{\linewidth}\raggedright
Холболт
\end{minipage} & \begin{minipage}[b]{\linewidth}\raggedright
Зорилтот мсж/с
\end{minipage} & \begin{minipage}[b]{\linewidth}\raggedright
EMQX холболт
\end{minipage} & \begin{minipage}[b]{\linewidth}\raggedright
EMQX хүлээн авсан мсж/с
\end{minipage} & \begin{minipage}[b]{\linewidth}\raggedright
Хаягдсан (\texttt{msg\_dropped})
\end{minipage} & \begin{minipage}[b]{\linewidth}\raggedright
CPU \%
\end{minipage} & \begin{minipage}[b]{\linewidth}\raggedright
Сул RAM (MiB)
\end{minipage} \\
\midrule
\endhead
\bottomrule
\endlastfoot
50 & 50 & & & & & \\
100 & 100 & & & & & \\
250 & 250 & & & & & \\
500 & 500 & & & & & \\
1000 & 1000 & & & & & \\
2000 & 2000 & & & & & \\
\end{longtable}\endgroup

\hypertarget{ux430ux43bux445ux430ux43c-5-emqx-rule-engine-ux43cux435ux441ux441ux435ux436ux438ux439ux433-ux431ux440ux43eux43aux435ux440-ux434ux44dux44dux440-ux447ux438ux433ux43bux4afux4afux43bux44dux445-30-ux43cux438ux43d-ux437ux43a}{%
\subsection{Алхам 5 --- EMQX rule engine: мессежийг брокер дээр чиглүүлэх (30 мин) {[}ЗК{]}}\label{ux430ux43bux445ux430ux43c-5-emqx-rule-engine-ux43cux435ux441ux441ux435ux436ux438ux439ux433-ux431ux440ux43eux43aux435ux440-ux434ux44dux44dux440-ux447ux438ux433ux43bux4afux4afux43bux44dux445-30-ux43cux438ux43d-ux437ux43a}}

\textbf{5.1} EMQX самбар (\texttt{:18083}) → \textbf{Integration → Rules → Create}. Чичиргээний хэт өндөр утгыг шүүх дүрэм:

\begin{Shaded}
\begin{Highlighting}[]
\KeywordTok{SELECT}
\NormalTok{  payload.device\_id     }\KeywordTok{as}\NormalTok{ device\_id,}
\NormalTok{  payload.vibration\_rms }\KeywordTok{as}\NormalTok{ vibration,}
\NormalTok{  payload.temperature   }\KeywordTok{as}\NormalTok{ temperature,}
\NormalTok{  clientid              }\KeywordTok{as} \KeywordTok{source}\NormalTok{,}
  \DataTypeTok{timestamp}             \KeywordTok{as}\NormalTok{ ts}
\KeywordTok{FROM}
  \OtherTok{"cnc302/+/+/+/+/telemetry"}
\KeywordTok{WHERE}
\NormalTok{  payload.vibration\_rms }\OperatorTok{\textgreater{}} \FloatTok{1.5}
\end{Highlighting}
\end{Shaded}

Үйлдэл (Action): \textbf{Republish} → сэдэв \texttt{cnc302/alerts/vibration}.

\textbf{5.2 Туршина:}

\begin{Shaded}
\begin{Highlighting}[]
\CommentTok{\# [ЗК] терминал 1 — дохиоллыг сонсоно; терминал 2 — гажилтай өгөгдөл илгээнэ}
\ExtensionTok{mosquitto\_sub} \AttributeTok{{-}h}\NormalTok{ localhost }\AttributeTok{{-}t} \StringTok{\textquotesingle{}cnc302/alerts/\#\textquotesingle{}} \AttributeTok{{-}v}
\ExtensionTok{python3}\NormalTok{ tools/sim\_device.py }\AttributeTok{{-}{-}target}\NormalTok{ cloud }\AttributeTok{{-}{-}host}\NormalTok{ localhost }\AttributeTok{{-}{-}devices}\NormalTok{ 5 }\AttributeTok{{-}{-}interval}\NormalTok{ 1 }\AttributeTok{{-}{-}anomaly{-}rate}\NormalTok{ 0.2}
\end{Highlighting}
\end{Shaded}

\textbf{5.3 Ачаалалтай үед дүрэм ажиллаж байна уу?} Алхам 4-ийн 500 холболтын шатыг дахин ажиллуулж, дүрмийн статистик (\texttt{Matched}, \texttt{Failed}, \texttt{Rate}) ба EMQX-ийн CPU-г дүрэмгүй үетэй харьцуул.

\textbf{5.4 Цонхны нэгтгэл.} EMQX CE-д доорх дэмжигдэх эсэхийг \textbf{шалгаж, үр дүнг бич} --- дэмжигдэхгүй бол энэ нь Лаб 5-д Node-RED хэрэгтэй болох шалтгаан:

\begin{Shaded}
\begin{Highlighting}[]
\KeywordTok{SELECT} \FunctionTok{avg}\NormalTok{(payload.temperature) }\KeywordTok{as}\NormalTok{ avg\_temp}
\KeywordTok{FROM} \OtherTok{"cnc302/+/+/+/+/telemetry"}
\KeywordTok{GROUP} \KeywordTok{BY}\NormalTok{ device\_id, }\DataTypeTok{INTERVAL} \DecValTok{60}\NormalTok{s}
\end{Highlighting}
\end{Shaded}

\hypertarget{ux445ux4afux441ux43dux44dux433ux442-4.6-ux431ux43eux43bux43eux432ux441ux440ux443ux443ux43bux430ux43bux442ux44bux43d-ux431ux430ux439ux440ux43bux430ux43b}{%
\subsubsection{Хүснэгт 4.6 --- Боловсруулалтын байрлал}\label{ux445ux4afux441ux43dux44dux433ux442-4.6-ux431ux43eux43bux43eux432ux441ux440ux443ux443ux43bux430ux43bux442ux44bux43d-ux431ux430ux439ux440ux43bux430ux43b}}

\begingroup\scriptsize\begin{longtable}[]{@{}
  >{\raggedright\arraybackslash}p{(\columnwidth - 8\tabcolsep) * \real{0.2000}}
  >{\raggedright\arraybackslash}p{(\columnwidth - 8\tabcolsep) * \real{0.2000}}
  >{\raggedright\arraybackslash}p{(\columnwidth - 8\tabcolsep) * \real{0.2000}}
  >{\raggedright\arraybackslash}p{(\columnwidth - 8\tabcolsep) * \real{0.2000}}
  >{\raggedright\arraybackslash}p{(\columnwidth - 8\tabcolsep) * \real{0.2000}}@{}}
\toprule
\begin{minipage}[b]{\linewidth}\raggedright
Байрлал
\end{minipage} & \begin{minipage}[b]{\linewidth}\raggedright
Дундаж саатал (мс)
\end{minipage} & \begin{minipage}[b]{\linewidth}\raggedright
Нэмэгдсэн CPU \%
\end{minipage} & \begin{minipage}[b]{\linewidth}\raggedright
500 холболт дээр ажиллав уу
\end{minipage} & \begin{minipage}[b]{\linewidth}\raggedright
Логикийн уян хатан (1--5)
\end{minipage} \\
\midrule
\endhead
\bottomrule
\endlastfoot
EMQX rule engine ({[}ЗК{]}) & & & & \\
Ирмэгийн агентын шүүлт ({[}Pi{]}, Лаб 1) & & & & \\
Node-RED ({[}ЗК{]}, Лаб 5-д хэмжинэ) & --- & --- & --- & \\
\end{longtable}\endgroup

Цонхны нэгтгэл EMQX CE-д дэмжигдэв үү: ule{1.5cm}{0.4pt} Хэрэв үгүй бол хаана хийх вэ: ule{1.5cm}{0.4pt}

\hypertarget{ux430ux43bux445ux430ux43c-6-ux445ux43eux451ux440-ux43cux443ux440ux443ux439-ux43dux44dux433-ux442ux44dux43dux445ux43bux44dux433-ux445ux430ux43dux430ux43bux442ux44bux43d-ux448ux438ux43dux436ux438ux43bux433ux44dux44d-30-ux43cux438ux43d-ux437ux43a}{%
\subsection{Алхам 6 --- Хоёр муруй, нэг тэнхлэг: ханалтын шинжилгээ (30 мин) {[}ЗК{]}}\label{ux430ux43bux445ux430ux43c-6-ux445ux43eux451ux440-ux43cux443ux440ux443ux439-ux43dux44dux433-ux442ux44dux43dux445ux43bux44dux433-ux445ux430ux43dux430ux43bux442ux44bux43d-ux448ux438ux43dux436ux438ux43bux433ux44dux44d-30-ux43cux438ux43d-ux437ux43a}}

\begin{Shaded}
\begin{Highlighting}[]
\CommentTok{\# [ЗК] хоёр хостын CSV{-}г нэг дор өгнө (Pi{-}гийнхийг эхлээд хуулж ав)}
\BuiltInTok{cd}\NormalTok{ \textasciitilde{}/cnc302}
\FunctionTok{scp}\NormalTok{ pi@}\OperatorTok{\textless{}}\NormalTok{PI{-}IP}\OperatorTok{\textgreater{}}\NormalTok{:\textasciitilde{}/cnc302/measurements/loadtest{-}edge{-}ramp{-}}\PreprocessorTok{*}\NormalTok{.csv measurements/}
\ExtensionTok{python3}\NormalTok{ lab04/plot\_scaling.py measurements/loadtest{-}edge{-}}\PreprocessorTok{*}\NormalTok{.csv }\DataTypeTok{\textbackslash{}}
\NormalTok{        measurements/loadtest{-}cloud{-}}\PreprocessorTok{*}\NormalTok{.csv }\AttributeTok{{-}{-}out}\NormalTok{ lab04/scaling.png}
\end{Highlighting}
\end{Shaded}

matplotlib байхгүй бол ASCII график гарна (edge = ●, cloud = ○) --- үзүүлэнд ч хангалттай.

\hypertarget{ux445ux4afux441ux43dux44dux433ux442-4.7-ux430ux43bux44c-ux43dux4e9ux4e9ux446-ux442ux4afux440ux4afux4afux43bux436-ux445ux430ux43dux430ux432}{%
\subsubsection{Хүснэгт 4.7 --- Аль нөөц түрүүлж ханав}\label{ux445ux4afux441ux43dux44dux433ux442-4.7-ux430ux43bux44c-ux43dux4e9ux4e9ux446-ux442ux4afux440ux4afux4afux43bux436-ux445ux430ux43dux430ux432}}

Багана бүрд \textbf{тоо ба эх сурвалж} бич. "RAM ханасан" гэсэн ганц үг оноо авахгүй.

\begingroup\scriptsize\begin{longtable}[]{@{}
  >{\raggedright\arraybackslash}p{(\columnwidth - 6\tabcolsep) * \real{0.2500}}
  >{\raggedright\arraybackslash}p{(\columnwidth - 6\tabcolsep) * \real{0.2500}}
  >{\raggedright\arraybackslash}p{(\columnwidth - 6\tabcolsep) * \real{0.2500}}
  >{\raggedright\arraybackslash}p{(\columnwidth - 6\tabcolsep) * \real{0.2500}}@{}}
\toprule
\begin{minipage}[b]{\linewidth}\raggedright
Нөөц
\end{minipage} & \begin{minipage}[b]{\linewidth}\raggedright
Ирмэг: ханасан уу, ямар тоогоор
\end{minipage} & \begin{minipage}[b]{\linewidth}\raggedright
Үүл: ханасан уу, ямар тоогоор
\end{minipage} & \begin{minipage}[b]{\linewidth}\raggedright
Нотолгоо (хүснэгт/багана)
\end{minipage} \\
\midrule
\endhead
\bottomrule
\endlastfoot
RAM (сул MiB, swap si/so) & & & \\
CPU (\% 400-аас) & & & \\
Уплинк (Mbit/с 95-аас) & & & \\
microSD (I/O хүлээлт \texttt{wa}) & & & \\
Холболтын тоо / файлын дескриптор & & & \\
\textbf{ХАМГИЙН ТҮРҮҮНД ханасан нь} & & & \\
\end{longtable}\endgroup

\hypertarget{ux445ux4afux441ux43dux44dux433ux442-4.8-ux442ux430ux430ux43cux430ux433-ux431ux430-ux431ux43eux434ux438ux442}{%
\subsubsection{Хүснэгт 4.8 --- Таамаг ба бодит}\label{ux445ux4afux441ux43dux44dux433ux442-4.8-ux442ux430ux430ux43cux430ux433-ux431ux430-ux431ux43eux434ux438ux442}}

\begingroup\scriptsize\begin{longtable}[]{@{}llll@{}}
\toprule
& Таамаг (Хүснэгт 4.2) & Бодит & Зөрүү \% \\
\midrule
\endhead
\bottomrule
\endlastfoot
Ирмэгийн дээд холболт & & & \\
Ирмэгийн дээд мсж/с & & & \\
Түрүүлж ханасан нөөц & & & таарсан уу: тийм/үгүй \\
Үүлний дээд холболт & & & \\
\end{longtable}\endgroup

Зөрүү 30\%-иас их бол шалтгааныг 2--4 өгүүлбэрээр бич. Боломжит шалтгаанууд: холболт бүрийн тогтмол зардал шугаман биш; дараалал дүүрэхэд саатал экспоненциалаар өснө; температур өсөж давтамж буурсан; swap эхэлсэн (microSD рүү бичих нь RAM-аас олон зуун дахин удаан).

\hypertarget{ux430ux43bux445ux430ux43c-7-ux431ux430ux433ux442ux430ux430ux43cux436ux438ux439ux43d-ux442ux43eux43cux44cux451ux43e-ux431ux430-ux442ux4e9ux441ux43bux438ux439ux43d-ux442ux43eux43eux446ux43eux43e-15-ux43cux438ux43d-ux437ux43a}{%
\subsection{Алхам 7 --- Багтаамжийн томьёо ба төслийн тооцоо (15 мин) {[}ЗК{]}}\label{ux430ux43bux445ux430ux43c-7-ux431ux430ux433ux442ux430ux430ux43cux436ux438ux439ux43d-ux442ux43eux43cux44cux451ux43e-ux431ux430-ux442ux4e9ux441ux43bux438ux439ux43d-ux442ux43eux43eux446ux43eux43e-15-ux43cux438ux43d-ux437ux43a}}

\begin{verbatim}
Шаардлагатай ирмэгийн зангилааны тоо =  ⌈ N × R ⌉ / C_edge × S

  N        — төхөөрөмжийн тоо
  R        — нэг төхөөрөмжийн мессежийн хурд (мсж/с)
  C_edge   — хэмжсэн багтаамж (мсж/с), 70% ашиглалтаар
  S        — аюулгүйн коэффициент (санал: 1.5)
\end{verbatim}

Өөрсдийн тоогоор бөглө:

\begin{verbatim}
C_edge (70%) = ______ мсж/с      C_cloud (70%) = ______ мсж/с
Төслийн ачаалал: ______ төхөөрөмж × ______ мсж/с = ______ мсж/с
→ Хэрэгтэй Pi: ______     → Нэг үүл хангалттай юу: ______
Нэг Pi-д ногдох дээд төхөөрөмж (RAM-аар хязгаарлавал): ______
\end{verbatim}

Дараа нь хариул: \textbf{ижил ачааллыг үүл дээр авах уу, эсвэл ирмэг нэмэх үү?} Уплинкийн зардал, тасалдалд тэсвэртэй байдал, төхөөрөмжийн үнэ гурвыг харьцуул (Лаб 5-д уплинк тасрахад юу болохыг үзнэ).

\hypertarget{ux430ux43bux445ux430ux43c-8-ux4afux437ux4afux4afux43bux44dux43dux433ux438ux439ux43d-ux431ux44dux43bux442ux433ux44dux43b-ux431ux430-commit-15-ux43cux438ux43d-ux437ux43a}{%
\subsection{Алхам 8 --- Үзүүлэнгийн бэлтгэл ба commit (15 мин) {[}ЗК{]}}\label{ux430ux43bux445ux430ux43c-8-ux4afux437ux4afux4afux43bux44dux43dux433ux438ux439ux43d-ux431ux44dux43bux442ux433ux44dux43b-ux431ux430-commit-15-ux43cux438ux43d-ux437ux43a}}

\textbf{10 минут, гурван гишүүн бүгд ярина:}

\begingroup\scriptsize\begin{longtable}[]{@{}
  >{\raggedright\arraybackslash}p{(\columnwidth - 4\tabcolsep) * \real{0.3333}}
  >{\raggedright\arraybackslash}p{(\columnwidth - 4\tabcolsep) * \real{0.3333}}
  >{\raggedright\arraybackslash}p{(\columnwidth - 4\tabcolsep) * \real{0.3333}}@{}}
\toprule
\begin{minipage}[b]{\linewidth}\raggedright
Хэсэг
\end{minipage} & \begin{minipage}[b]{\linewidth}\raggedright
Хугацаа
\end{minipage} & \begin{minipage}[b]{\linewidth}\raggedright
Хэн
\end{minipage} \\
\midrule
\endhead
\bottomrule
\endlastfoot
1. Таамаг ба арифметик (Хүснэгт 4.2), аргачлал & 3 мин & гишүүн 1 \\
2. Хоёр муруй нэг тэнхлэгт, аль нөөц түрүүлж ханав (4.3--4.5, 4.7) & 4 мин & гишүүн 2 \\
3. Rule engine, багтаамжийн томьёо, төслийн тооцоо (4.6, 4.8) & 3 мин & гишүүн 3 \\
\end{longtable}\endgroup

\textbf{Заавал үзүүлэх:} ханалтын график, хязгаарын цэг, таамаг--бодит зөрүү, ханасан нөөцийн \textbf{нотолгоо}.

\begin{Shaded}
\begin{Highlighting}[]
\BuiltInTok{cd}\NormalTok{ \textasciitilde{}/cnc302}
\FunctionTok{git}\NormalTok{ add lab04/}
\CommentTok{\# measurements/ нь .gitignore{-}д — тайланд хэрэгтэй CSV{-}г зориудаар албадан нэмнэ}
\FunctionTok{git}\NormalTok{ add }\AttributeTok{{-}f}\NormalTok{ measurements/loadtest{-}edge{-}ramp{-}}\PreprocessorTok{*}\NormalTok{.csv measurements/loadtest{-}cloud{-}ramp{-}}\PreprocessorTok{*}\NormalTok{.csv}
\FunctionTok{git}\NormalTok{ status }\AttributeTok{{-}{-}short}
\FunctionTok{git}\NormalTok{ commit }\AttributeTok{{-}m} \StringTok{"Лаб 4: ирмэг ба үүлний багтаамжийн хязгаар"}
\FunctionTok{git}\NormalTok{ tag lab04{-}done }\KeywordTok{\&\&} \FunctionTok{git}\NormalTok{ push }\KeywordTok{\&\&} \FunctionTok{git}\NormalTok{ push }\AttributeTok{{-}{-}tags}
\CommentTok{\# Нууц файл ороогүйг ЗААВАЛ шалга}
\FunctionTok{git}\NormalTok{ ls{-}files }\KeywordTok{|} \FunctionTok{grep} \AttributeTok{{-}E} \StringTok{\textquotesingle{}\textbackslash{}.env$|bridge\textbackslash{}.conf$|\textbackslash{}.key$|\textbackslash{}.crt$|devices\textbackslash{}.csv\textquotesingle{}}   \CommentTok{\# хоосон байх ЁСТОЙ}
\end{Highlighting}
\end{Shaded}

\hypertarget{ux445ux44fux43dux430ux43bux442ux44bux43d-ux430ux441ux443ux443ux43bux442ux443ux443ux434-ux4afux437ux4afux4afux43bux44dux43dux434-ux430ux43cux430ux430ux440-ux445ux430ux440ux438ux443ux43bux43dux430}{%
\section{5. Хяналтын асуултууд (үзүүлэнд амаар хариулна)}\label{ux445ux44fux43dux430ux43bux442ux44bux43d-ux430ux441ux443ux443ux43bux442ux443ux443ux434-ux4afux437ux4afux4afux43bux44dux43dux434-ux430ux43cux430ux430ux440-ux445ux430ux440ux438ux443ux43bux43dux430}}

\begin{enumerate}
\def\labelenumi{\arabic{enumi}.}
\item
  Хүснэгт 4.7-д ирмэг дээр аль нөөц түрүүлж ханав? Тухайн шатанд сул RAM хэдэн MiB, CPU хэдэн \%, уплинк хэдэн Mbit/с байв --- гурвуулангийн тоог хэл. Яагаад нөгөө хоёр нь хязгаарлагч биш вэ?
\item
  Хүснэгт 4.2-ын таамаг ба Хүснэгт 4.8-ын бодит үр дүн хэдэн хувиар зөрөв? Шугаман экстраполяци яг хаанаас алдагдав?
\item
  Хүснэгт 4.3-т swap \texttt{si/so} тэгээс өссөн шат аль вэ? Түүнээс хойшхи саатлын тоог яагаад тайланд оруулж болохгүй вэ?
\item
  Хүснэгт 4.4-т тооцоолсон ба хэмжсэн Mbit/с хэр таарав? Зөрүү гарсан бол хаана алдагдав (TCP толгой, MQTT толгой, ACK урсгал, гүүрний QoS)?
\item
  Хүснэгт 4.3 ба 4.5-ын дээд мсж/с хэдэн дахин зөрөв? Тэр харьцаа нь техник хангамжийн (RAM, цөм) харьцаатай таарч байна уу, эсвэл програм хангамжийн (mosquitto vs EMQX) ялгаа нь давамгайлав уу?
\item
  Хүснэгт 4.6-д rule engine 500 холболт дээр ажиллав уу? Хэрэв EMQX-ийн CPU дүрэмтэй үед мэдэгдэхүйц өссөн бол 10 000 төхөөрөмжийн флотод юу хийх вэ?
\item
  Таны төслийн тооцоогоор (Алхам 7) хэдэн Pi хэрэгтэй гарав? Хэрэв ирмэг дээр 90\% шүүлт хийж, зөвхөн 10\%-ийг үүл рүү илгээвэл энэ тоо хэрхэн өөрчлөгдөх вэ --- уплинкийн арифметикээр хариул.
\end{enumerate}

\hypertarget{ux445ux4afux43bux44dux44dux43bux433ux44dux43d-ux4e9ux433ux4e9ux445-ux437ux4afux439ux43b}{%
\section{6. Хүлээлгэн өгөх зүйл}\label{ux445ux4afux43bux44dux44dux43bux433ux44dux43d-ux4e9ux433ux4e9ux445-ux437ux4afux439ux43b}}

\begingroup\scriptsize\begin{longtable}[]{@{}
  >{\raggedright\arraybackslash}p{(\columnwidth - 4\tabcolsep) * \real{0.3333}}
  >{\raggedright\arraybackslash}p{(\columnwidth - 4\tabcolsep) * \real{0.3333}}
  >{\raggedright\arraybackslash}p{(\columnwidth - 4\tabcolsep) * \real{0.3333}}@{}}
\toprule
\begin{minipage}[b]{\linewidth}\raggedright
\#
\end{minipage} & \begin{minipage}[b]{\linewidth}\raggedright
Зүйл
\end{minipage} & \begin{minipage}[b]{\linewidth}\raggedright
Байрлал
\end{minipage} \\
\midrule
\endhead
\bottomrule
\endlastfoot
1 & 10 минутын үзүүлэн (гурван гишүүн) & хичээл дээр \\
2 & Ачааллын муруй (хоёр зорилт нэг тэнхлэгт) & \texttt{lab04/scaling.png} эсвэл ASCII гаралт \\
3 & Хоёр хостын түүхий CSV & \texttt{measurements/\allowbreak{}loadtest-edge-*.\allowbreak{}csv}, \texttt{loadtest-cloud-*.csv} \\
4 & Хүснэгт 4.1--4.8 (слайд эсвэл нэг хуудас) & \texttt{lab04/tables.md} \\
5 & Git tag & \texttt{lab04-done} \\
\end{longtable}\endgroup

\hypertarget{ux4afux43dux44dux43bux433ux44dux44dux43dux438ux439-ux448ux430ux43bux433ux443ux443ux440-10-ux43eux43dux43eux43e}{%
\section{7. Үнэлгээний шалгуур (10 оноо)}\label{ux4afux43dux44dux43bux433ux44dux44dux43dux438ux439-ux448ux430ux43bux433ux443ux443ux440-10-ux43eux43dux43eux43e}}

\begingroup\scriptsize\begin{longtable}[]{@{}
  >{\raggedright\arraybackslash}p{(\columnwidth - 2\tabcolsep) * \real{0.5000}}
  >{\raggedright\arraybackslash}p{(\columnwidth - 2\tabcolsep) * \real{0.5000}}@{}}
\toprule
\begin{minipage}[b]{\linewidth}\raggedright
Шалгуур
\end{minipage} & \begin{minipage}[b]{\linewidth}\raggedright
Оноо
\end{minipage} \\
\midrule
\endhead
\bottomrule
\endlastfoot
Хоёр зорилт (ирмэг ба үүл) бүрэн шатаар ачаалагдаж, хязгаар нь олдсон & 3 \\
\textbf{Түрүүлж ханасан нөөц зөв тогтоогдож, тоогоор нотлогдсон} (Хүснэгт 4.7) & 3 \\
Уплинкийн арифметик хийгдэж, хэмжилттэй харьцуулагдсан (4.2, 4.4) & 2 \\
Rule engine ажиллаж, ачаалалтай үед хэмжигдсэн & 1 \\
Гурван гишүүн бүгд ярьж, асуултад хариулсан & 1 \\
\end{longtable}\endgroup

\textbf{Онцгой оноо (+1):} ханалтын муруйн \textbf{буурах} хэсгийг баримтжуулсан бол (ачаалал нэмэгдэхэд хүлээн авалт буурсан).

\textbf{Оноо хасах гол шалтгаан:} ханасан нөөцийг нотолгоогүй нэрлэх; throttle \texttt{0x0} биш үеийн тоог хүчинтэй мэт үзүүлэх.

\hypertarget{ux442ux4afux433ux44dux44dux43cux44dux43b-ux430ux43bux434ux430ux430}{%
\section{8. Түгээмэл алдаа}\label{ux442ux4afux433ux44dux44dux43cux44dux43b-ux430ux43bux434ux430ux430}}

\begingroup\scriptsize\begin{longtable}[]{@{}
  >{\raggedright\arraybackslash}p{(\columnwidth - 4\tabcolsep) * \real{0.3333}}
  >{\raggedright\arraybackslash}p{(\columnwidth - 4\tabcolsep) * \real{0.3333}}
  >{\raggedright\arraybackslash}p{(\columnwidth - 4\tabcolsep) * \real{0.3333}}@{}}
\toprule
\begin{minipage}[b]{\linewidth}\raggedright
Шинж
\end{minipage} & \begin{minipage}[b]{\linewidth}\raggedright
Шалтгаан
\end{minipage} & \begin{minipage}[b]{\linewidth}\raggedright
Шийдэл
\end{minipage} \\
\midrule
\endhead
\bottomrule
\endlastfoot
emqtt-bench 1024 холболт дээр зогсоно & \texttt{ulimit\ -n} бага & \texttt{ulimit\ -n\ 65535} \\
Ачаалал өгсөн ч Pi дээр юу ч өөрчлөгдөхгүй & ачаалал үүлний EMQX рүү явж байна & \texttt{-\/-target\ edge} ба \texttt{\textless{}PI-IP\textgreater{}}-ээ шалга \\
Уплинк огт ачаалагдахгүй & сэдэв \texttt{cnc302/\textless{}SITE\textgreater{}/} угтваргүй & Алхам 3.3-ын \texttt{TOPIC=} \\
\texttt{collect\_metrics.py} EMQX API алдаа & нууц үг буруу & \texttt{export\ EMQX\_\allowbreak{}DASHBOARD\_\allowbreak{}PASSWORD=\allowbreak{}…} эсвэл \texttt{-\/-no-emqx} \\
CSV-д \texttt{emqx\_connections} ирмэгийн ачааллыг харуулахгүй & тэр багана үүлний EMQX-ийнх & \texttt{\$SYS/\allowbreak{}broker/\allowbreak{}clients/\allowbreak{}connected}-оос гараар бич \\
Pi гэнэт хариу өгөхөө болино & swap + microSD дүүрсэн & тэжээл унтраахаас өмнө 2 мин хүлээ; шатыг буулга \\
CPU 400\% харагдана & 4 цөмийн нийлбэр & 100\% = нэг цөм, 400\% = бүрэн ханалт \\
\texttt{throttled} ≠ \texttt{0x0} & халалт эсвэл сул тэжээл & хөргөлт, 5 V/2.5 A тэжээл; хэмжилтийг \textbf{дахин} хий \\
Компьютер өөрөө эхэлж ханана & ачаалал үүсгэгч нэг машин дээр & өөр машинаас ачаал, эсвэл \texttt{-c} багасгаж хурдыг нэмэгдүүл \\
\end{longtable}\endgroup
